% Detection vs Forecasting Comparison Subsection
\subsection{Detection Versus Forecasting}

Detection and forecasting serve different clinical purposes. Detection aims to identify seizures as they occur for timely intervention. Forecasting aims to predict seizures before onset to allow preventive action. These different objectives lead to distinct technical approaches and validation requirements.

\subsubsection{Clinical Objectives}

Detection systems alert caregivers when a seizure occurs. This enables intervention during or immediately after the seizure. Detection is most valuable for generalized tonic-clonic or focal to bilateral tonic-clonic seizures where caregiver response can prevent injury or complications.

Forecasting systems provide advance warning of increased seizure likelihood. This allows patients to modify activities, take rescue medication, or avoid risky situations. Forecasting could improve quality of life by reducing the uncertainty of when seizures might occur.

\subsubsection{Algorithmic Differences}

Detection and forecasting employ different algorithmic paradigms. Detection research uses diverse approaches including feature-based methods (3 studies), ensemble CNN (1 study), two-stage architectures (1 study), and anomaly detection (2 studies). Forecasting research concentrates on LSTM-based approaches (3 studies) with some use of hybrid architectures.

Detection architectures prioritize spatial pattern recognition. CNN-based approaches appear in two detection studies but no forecasting studies. Forecasting architectures prioritize temporal dependency modeling. LSTM or LSTM-ensemble approaches appear in three of five forecasting studies.

\subsubsection{Modality Preferences}

Detection and forecasting favor different biosignal modalities. Detection studies emphasize movement sensors. Six of eight detection studies use accelerometer data. Three use gyroscope data. Movement sensors capture the motor manifestations of convulsive seizures.

Forecasting studies emphasize autonomic nervous system indicators. Four of five forecasting studies use EDA. Three use ECG or HRV data. Two use temperature data. Autonomic changes may precede seizure onset by minutes to hours.

\subsubsection{Performance Comparison}

Detection and forecasting cannot be directly compared because they use different metrics. Detection sensitivity ranges from 38\% to 100\%. Forecasting performance measured by AUC ranges from 0.74 to 0.75. These metrics assess different constructs.

Detection faces the FPR challenge. Only two of eight detection studies meet the clinical FPR benchmark of below 0.05 to 0.1 per hour. High FPR leads to alarm fatigue and limits real-world utility. Forecasting does not have an established clinical benchmark for acceptable performance.

\subsubsection{Clinical Maturity}

Detection shows greater clinical maturity than forecasting. Two detection studies meet Phase 3 FPR benchmarks. Commercial devices exist for detection including the Embrace watch and NightWatch armband. No forecasting study meets a clear clinical benchmark. No device has regulatory approval for seizure forecasting.

Detection benefits from clearer clinical requirements. The ILAE has defined guidelines for detection device validation including FPR targets. Forecasting lacks equivalent standardized guidelines. This gap complicates assessment of forecasting readiness.

\subsubsection{Sample Size Disparity}

Detection studies enroll larger samples than forecasting studies. Detection studies have a median sample size of 44 participants. Forecasting studies have a median sample size of 11 participants. This disparity may reflect different validation requirements.

Larger detection studies enable more robust performance estimates. \textcite{spahrDeepLearningbasedDetection2025} enrolled 384 patients across eight centers. \textcite{dongTwoLayerEnsembleMethod2022} enrolled 68 patients with home monitoring. The largest forecasting study, \textcite{vielufSeizureMonitoringCombined2025}, enrolled 70 patients.

\subsubsection{Validation Setting}

Detection and forecasting differ in validation settings. Six of eight detection studies were conducted in hospital or EMU settings. Only three detection studies included home validation. Three of five forecasting studies were conducted in home settings.

Home validation may be particularly important for forecasting. Forecasting systems must learn from normal daily activities to avoid false alarms from exercise, stress, or other confounding factors. Hospital settings may not capture this diversity of activities.

\textcite{nasseriAmbulatorySeizureForecasting2021} achieved the most rigorous validation with concurrent invasive EEG confirmation in ambulatory settings. However, only six patients were studied.

\subsubsection{Reporting Practices}

Detection and forecasting differ in metric reporting. Forecasting studies consistently report patient-level success rates. Detection studies rarely report individual patient outcomes. This reporting gap complicates assessment of which patients benefit from each approach.

Forecasting studies report forecasting-specific metrics including IoC and TiW. Detection studies emphasize sensitivity and FPR. These different metrics reflect different clinical priorities but complicate cross-domain comparison.

\textbf{Key finding.} Detection is closer to clinical translation than forecasting. Two detection approaches meet clinical FPR benchmarks. Forecasting shows consistent AUC around 0.75 but lacks clear clinical benchmarks and regulatory pathways.

\subsubsection{Performance Ceiling}

Forecasting shows a consistent performance ceiling across studies. Three forecasting studies report AUC between 0.74 and 0.75. This consistency across different methods and modalities suggests a fundamental limit of non-invasive forecasting approaches.

Detection has reached a sensitivity ceiling of 100\% but struggles with FPR. The best performing detection study, \textcite{spahrDeepLearningbasedDetection2025}, achieves 96\% sensitivity with 0.005/h FPR. However, most detection studies report FPR above clinical acceptability thresholds.

\textbf{Limitation identified.} Direct comparison between detection and forecasting is challenging due to different metrics, validation approaches, and clinical objectives. The field lacks standardized frameworks for assessing relative clinical value.
